\documentclass[10pt]{article}
\usepackage[margin=0.85in]{geometry}
\usepackage{amsmath,amssymb,amsthm,hyperref}
\newtheorem{theorem}{Theorem}
\title{A counterexample to the Riordan nullity conjecture}
\author{Principia Math in Progress}
\date{13 September 2026}
\begin{document}
\maketitle

\paragraph{Task 1. Resolve the nullity equality conjecture.}
Conjecture 7.9 of Cheon, Jung, Kitaev and Mojallal \cite{CJKM} asserts
that every singular io-decomposable Riordan graph of the Bell type has
\(\eta(A_n)=\eta(B)\), where, after listing odd vertices first,
\[
 A_n\sim\begin{pmatrix} A_{\lceil n/2\rceil}&B\\B^T&0\end{pmatrix}.
\]
Here nullity is the dimension of the right kernel. In particular,
\(B\) has \(\lfloor n/2\rfloor\) columns. We give a counterexample
with \(n=15\), \(\eta(A_{15})=1\), and \(\eta(B)=0\).

For clarity, the Bell-type graph \(G_n(g,zg)\) has vertices
\(1,\ldots,n\), zero diagonal, and, for \(i>j\), adjacency entry
\begin{equation}\label{entry}
 a_{ij}=a_{ji}=[z^{i-j-1}]g(z)^j\pmod2\quad\hbox{in }\{0,1\}.
\end{equation}
It is io-decomposable when its even vertices form an independent
set and its odd induced subgraph, relabelled in order, is
\(G_{\lceil n/2\rceil}(g,zg)\).

\begin{theorem}
Let
\[
 E(z)=1+z^6+z^8+z^{12},\qquad
 g(z)=\sum_{s=0}^{\infty}z^{2^s-1}E(z^{2^s}).
\]
Then \(G_{15}(g,zg)\) is a singular, proper, io-decomposable
Riordan graph of the Bell type with
\(\eta(A_{15})=1\ne0=\eta(B)\).
\end{theorem}

\begin{proof}
The displayed series is a well-defined element of \(\mathbb Z[[z]]\)
with \(g(0)=1\). Thus the Riordan pair is proper. It satisfies
\(g(z)=E(z)+zg(z^2)\). Working modulo 2, its odd part is
\(zg(z^2)\), and \(g(z)^2=g(z^2)\). For \(i>j\), the
even-even adjacency entry is
\[
 [z^{2(i-j)-1}]g(z)^{2j}=0.
\]
The odd part of \(g(z)^{2j-1}=g(z)(g(z)^{j-1})^2\) is
\(zg(z^2)^j\). Consequently
\[
 a_{2i-1,2j-1}=[z^{2(i-j)-1}]g(z)^{2j-1}
              =[z^{i-j-1}]g(z)^j=a_{ij}\pmod2.
\]
This proves io-decomposability for every order, including 15.

Only the following initial coefficients are needed for \(A_{15}\):
\[
 g(z)\equiv1+z+z^3+z^6+z^7+z^8+z^{12}+z^{13}
                  \pmod {z^{14}}.
\]
Using \eqref{entry}, in the vertex order
\((1,3,5,7,9,11,13,15;2,4,6,8,10,12,14)\) its blocks are
\[
 C=A_8=\begin{pmatrix}
0&1&1&0&1&0&0&1\\
1&0&1&0&1&0&0&0\\
1&1&0&1&1&1&0&0\\
0&0&1&0&1&0&0&0\\
1&1&1&1&0&1&1&0\\
0&0&1&0&1&0&1&0\\
0&0&0&0&1&1&0&1\\
1&0&0&0&0&0&1&0
\end{pmatrix},\qquad
 B=\begin{pmatrix}
1&0&0&1&1&0&1\\
1&1&1&0&0&0&1\\
1&1&1&0&1&1&1\\
0&0&1&1&1&1&1\\
1&1&1&1&1&0&0\\
0&0&1&0&1&1&1\\
0&0&0&0&1&1&1\\
1&0&0&0&0&0&1
\end{pmatrix}.
\]
Let
\[
 o=(1,0,1,0,-1,0,-1,-1)^T,\qquad
 e=(-1,0,-1,-1,2,-1,1)^T.
\]
Direct integer multiplication gives
\[
 B^To=0,\qquad
 Co=(-1,1,0,0,1,-1,-2,0)^T=-Be.
\]
Thus \((o,e)^T\ne0\) is in the kernel of the reordered adjacency
matrix. In the original vertex order it is
\[
 (1,-1,0,0,1,-1,0,-1,-1,2,0,-1,-1,1,-1)^T.
\]

We also give a short exact certificate that \(B\) has full column
rank. Write \(b_i\) for row \(i\) of \(B\), and write
\(e_j^T\) for the standard row vectors in \(\mathbb R^7\).
Successively, the rows \(b_2,\ldots,b_8\) span
\begin{align*}
 e_4^T&=b_4-b_6,&e_3^T&=b_6-b_7,&e_2^T&=b_2-b_8-e_3^T,\\
 e_7^T&=b_7-b_3+b_2,&e_1^T&=b_8-e_7^T,\\
 e_5^T&=b_5-e_1^T-e_2^T-e_3^T-e_4^T,&
 e_6^T&=b_3-b_2-e_5^T.
\end{align*}
Hence \(\operatorname{rank}B=7\) and \(\eta(B)=0\).
As an additional check, deleting the first row of \(B\) gives
a square matrix of determinant \(-1\).

Finally, projection of \(\ker A_{15}\) onto its odd coordinates
is injective: if those coordinates vanish, the other block equation
is \(Be=0\), which forces \(e=0\). Its image lies in
\(\ker B^T\), which has dimension \(8-7=1\). The exhibited
nonzero kernel vector therefore proves \(\eta(A_{15})=1\).
All hypotheses of Conjecture 7.9 hold and its equality fails.
\end{proof}

\paragraph{Scope.}
This refutes the full universal conjecture by a concrete graph in the
intended class. The series is specified to all orders; an arbitrary
finite coefficient list is not assumed to extend to that class.
The example also has nonzero odd coordinates in a zero eigenvector,
contrary to the equivalent Conjecture 7.7. That equivalent formulation
receives no additional campaign credit.

\paragraph{Task 2. Classify counterexamples of order fifteen.}
Determine all generating-series prefixes giving counterexamples of
order fifteen and classify the resulting graphs up to isomorphism.

\paragraph{Task 3. Settle the even-order restriction.}
Determine whether \(\eta(A_n)=\eta(B)\) holds for every even-order
io-decomposable Bell-type Riordan graph, or exhibit an even-order
counterexample.

\paragraph{Task 4. Characterize families with nullity equality.}
Use the restriction of the odd adjacency quadratic form to
\(\ker B^T\) to identify conditions on \(g\) that force nullity
equality at every order.

\paragraph{Task 5. Construct an infinite family.}
Find an explicit family of distinct orders with counterexamples,
tracking how the nonzero odd kernel coordinates propagate under
the recursive odd-vertex decomposition.

\begin{thebibliography}{1}
\bibitem{CJKM} G.-S.~Cheon, J.-H.~Jung, S.~Kitaev and S.~A.~Mojallal,
\emph{Riordan graphs II: Spectral properties},
\href{https://arxiv.org/abs/1801.07021}{arXiv:1801.07021};
\emph{Linear Algebra and its Applications} \textbf{575} (2019),
174--215, \href{https://doi.org/10.1016/j.laa.2019.04.011}
{doi:10.1016/j.laa.2019.04.011}.
Definitions 2.2 and 2.9, Theorem 7.5, and Conjecture 7.9
appear on pp.~3, 6, 27, and 29 of the inspected 33-page
arXiv and accepted author manuscripts.
\end{thebibliography}
\end{document}
